Our findings reveal that introducing a judgmental tone to LLM prompts acts as a form of ``social facilitation,'' compelling the model to engage in deeper, more rigorous reasoning. Unlike the ``Sycophancy Trap'' observed in prior literature where models capitulate to explicit misinformation or specific false hints \cite{chang2026diagnosing, kim2025challenging}, our prompts challenged the general \textit{quality} of the reasoning. In response to being broadly ``judged,'' the model elaborated extensively, self-corrected on difficult social nuances, and avoided arbitrary answer flipping. 

\para{The Sweet Spot of Skepticism}
The results suggest the existence of a constructive prompting threshold. For subjective and highly contextual domains such as social arbitration (\raio), demanding that the model ``prove its reasoning can withstand intense scrutiny'' functions effectively as a zero-shot performance booster. By anticipating adversarial review, the LLM expands its search space for edge cases, yielding judgments that are significantly more aligned with human consensus. 

\para{Limitations}
Our study has several limitations that warrant consideration. First, we observed a strict ceiling effect on the \causalt dataset, where \gptmini achieved 100\% accuracy across all conditions. This prevented us from analyzing whether the Judgment Effect provides tangible benefits for strict logical and causal reasoning, or if the lengthened reasoning traces merely reflect redundant verbosity. Second, our pilot experiments were conducted exclusively on \gptmini with a sample size of 30 cases per dataset. It remains entirely plausible that different model families---particularly larger models like GPT-4o or heavily safety-aligned models like Claude---might exhibit a different ``apology threshold'' where the pressure transitions from constructive facilitation to destructive sycophancy.

\para{Broader Implications}
These findings highlight the profound impact of conversational framing on AI behavior. When deploying LLM-as-a-Judge systems in legal, ethical, or moderation contexts, system designers should carefully consider the persona and tone adopted in their meta-prompts. Purposefully injecting skepticism or warning against ``lazy'' thinking may serve as a simple, effective mechanism to maximize model utility and rigor without requiring complex algorithmic interventions.